\subsection{AI Service Implementation}

The AI service is implemented as a dedicated backend module responsible for communicating with the Groq cloud platform.\\

The system first assembles verified certificate information into a structured context object. This context is then combined with the user's question and submitted to the language model.

\begin{lstlisting}[
        language=typescript,
        style=codeStyle
    ]
    const context = `
      Verified Certificate Information

      Student Name: ${certificateData.blockchain.studentName}

      Token ID: ${certificateData.blockchain.tokenId}

      Degree: ${certificateData.blockchain.degree}

      Major: ${certificateData.blockchain.major}

      Wallet: ${certificateData.database.studentWallet}

      IPFS Hash: ${certificateData.blockchain.ipfsHash}

      Status: ${certificateData.database.lifecycleStatus}
    `;
\end{lstlisting}

The context is subsequently injected into the LLM request:

\begin{lstlisting}[
        language=typescript,
        style=codeStyle
    ]
    return await groq.chat.completions.create({
      model: "llama-3.3-70b-versatile",

      messages: [
        {
          role: "system",
          content: prompt,
        },
        {
          role: "user",
          content: `
            Question:
            ${question}

            Certificate Payload Context:
            ${context}
            `,
        },
      ],

      temperature: 0.1,
      stream: true,
    });
\end{lstlisting}

The AI-generated response is streamed through the Express backend and rendered progressively within the React chat interface.